\documentclass[11pt,letterpaper]{article}
\newcommand{\labid}{1-G}
\newcommand{\labtitleshort}{Receiver Firmware}
% Shared preamble for Module 1 lab handouts.
% Usage: each handout defines \labid and \labtitleshort, then % Shared preamble for Module 1 lab handouts.
% Usage: each handout defines \labid and \labtitleshort, then % Shared preamble for Module 1 lab handouts.
% Usage: each handout defines \labid and \labtitleshort, then \input{LabHandout_Preamble}
% BEFORE \begin{document}. Compiles with a single pdflatex pass (no glossaries tool).

\usepackage[T1]{fontenc}
\usepackage[margin=1in]{geometry}
\usepackage{titlesec}
\usepackage{enumitem}
\usepackage{booktabs}
\usepackage{tabularx}
\usepackage{graphicx}
\graphicspath{{Images/}{../Images/}}
\usepackage{xcolor}
\usepackage{hyperref}
\usepackage{fancyhdr}
\usepackage{amsmath}
\usepackage{amssymb}
\usepackage{tcolorbox}
\usepackage{caption}
\tcbuselibrary{skins,breakable}

% ── Glossary (per-document, built from shared glossary.tex) ──
% Uses \makenoidxglossaries so no external makeglossaries tool is
% needed: sorting and printing is done by LaTeX itself across two
% pdflatex passes (same cadence as the TOC). Only acronyms that
% are actually \gls{}-used in the handout appear in the printed
% glossary, and each \gls{} reference is a live hyperlink that
% jumps to its entry in the Acronyms section at the end.
\usepackage[acronym,nomain,nopostdot]{glossaries}
\makenoidxglossaries
\input{../../glossary}
% Call \printhandoutglossary just before \end{document} in each handout.
\newcommand{\printhandoutglossary}{%
	\clearpage
	\printnoidxglossary[type=\acronymtype,title={Glossary}]%
}

% ── Colors (match Module 1 lesson plan) ──────────────────────
\definecolor{stevens}{RGB}{163,31,52}
\definecolor{darkgray}{RGB}{60,60,60}
\definecolor{lightgray}{RGB}{240,240,240}
\definecolor{accentblue}{RGB}{30,80,140}

\hypersetup{
	colorlinks=true,
	linkcolor=stevens,
	urlcolor=accentblue,
	citecolor=darkgray,
}

% ── Defaults for header variables (handout overrides these) ──
\providecommand{\labid}{?}
\providecommand{\labtitleshort}{Lab Handout}

% ── Running header / footer ──────────────────────────────────
\pagestyle{fancy}
\fancyhf{}
\fancyhead[L]{\small\textcolor{darkgray}{Module 1 --- Lab Handout \labid}}
\fancyhead[R]{\small\textcolor{darkgray}{\labtitleshort}}
\fancyfoot[C]{\small\thepage}
\renewcommand{\headrulewidth}{0.4pt}

% ── Section formatting ───────────────────────────────────────
\titleformat{\section}{\Large\bfseries\color{stevens}}{\thesection}{0.6em}{}
\titleformat{\subsection}{\large\bfseries\color{darkgray}}{\thesubsection}{0.6em}{}

% ── Box styles ───────────────────────────────────────────────
\tcbset{
	infobox/.style={
		enhanced, colback=blue!3, colframe=accentblue, fonttitle=\bfseries,
		title={#1}, coltitle=white, breakable, left=6pt, right=6pt, top=4pt, bottom=4pt
	},
	warning/.style={
		enhanced, colback=red!3, colframe=stevens, fonttitle=\bfseries,
		title={#1}, coltitle=white, breakable, left=6pt, right=6pt, top=4pt, bottom=4pt
	},
	deliverable/.style={
		enhanced, colback=gray!5, colframe=darkgray, fonttitle=\bfseries,
		title={#1}, coltitle=white, breakable, left=6pt, right=6pt, top=4pt, bottom=4pt
	},
}

% ── Figure helper: always draws a placeholder box (images suppressed) ─
% Usage: \labfigure{filename}{width}{caption}{label}
\newcommand{\labfigure}[4]{%
	\begin{figure}[h!]
		\centering
		\fbox{\begin{minipage}[c][2.2in][c]{#2}
			\centering\color{darkgray}\small
				\textbf{[Figure placeholder]}\\[4pt]
				\texttt{\detokenize{#1}}\\[4pt]
				\textit{#3}
		\end{minipage}}%
		\caption{#3}
		\label{#4}
	\end{figure}%
}

% ── Title banner (call at top of document body) ──────────────
\newcommand{\labheader}[2]{%
	\begin{center}
		{\LARGE\bfseries\color{stevens} Lab Handout #1}\\[3pt]
		{\Large\bfseries #2}\\[6pt]
		{\small\color{darkgray} Capstone --- Module 1: \RF{} Hardware Foundations}
	\end{center}
	\vspace{4pt}
	{\color{darkgray}\rule{\textwidth}{0.4pt}}
	\vspace{10pt}
}

% BEFORE \begin{document}. Compiles with a single pdflatex pass (no glossaries tool).

\usepackage[T1]{fontenc}
\usepackage[margin=1in]{geometry}
\usepackage{titlesec}
\usepackage{enumitem}
\usepackage{booktabs}
\usepackage{tabularx}
\usepackage{graphicx}
\graphicspath{{Images/}{../Images/}}
\usepackage{xcolor}
\usepackage{hyperref}
\usepackage{fancyhdr}
\usepackage{amsmath}
\usepackage{amssymb}
\usepackage{tcolorbox}
\usepackage{caption}
\tcbuselibrary{skins,breakable}

% ── Glossary (per-document, built from shared glossary.tex) ──
% Uses \makenoidxglossaries so no external makeglossaries tool is
% needed: sorting and printing is done by LaTeX itself across two
% pdflatex passes (same cadence as the TOC). Only acronyms that
% are actually \gls{}-used in the handout appear in the printed
% glossary, and each \gls{} reference is a live hyperlink that
% jumps to its entry in the Acronyms section at the end.
\usepackage[acronym,nomain,nopostdot]{glossaries}
\makenoidxglossaries
% Shared acronym macro container for EE800/820 lesson plan modules.
% Usage: type \IDE, \FPGA, \UART, etc. First use is expanded; later uses show acronym only.
% Requires in each document preamble:
%   \usepackage[acronym,toc]{glossaries}
%   \makeglossaries
%   \input{../glossary}
% Print used acronyms near end of document:
%   \printglossary[type=\acronymtype,title={Acronyms Used}]

% Acronym definitions (only used entries appear in the printed glossary).
\newacronym{ask}{ASK}{Amplitude Shift Keying}
\newacronym{bga}{BGA}{Ball Grid Array}
\newacronym{bom}{BOM}{Bill of Materials}
\newacronym{bram}{BRAM}{Block Random Access Memory}
\newacronym{crc}{CRC}{Cyclic Redundancy Check}
\newacronym{css}{CSS}{Chirp Spread Spectrum}
\newacronym{dma}{DMA}{Direct Memory Access}
\newacronym{dsp}{DSP}{Digital Signal Processing}
\newacronym{eda}{EDA}{Exploratory Data Analysis}
\newacronym{fifo}{FIFO}{First In, First Out}
\newacronym{fpga}{FPGA}{Field-Programmable Gate Array}
\newacronym{fsk}{FSK}{Frequency Shift Keying}
\newacronym{fsm}{FSM}{Finite State Machine}
\newacronym{gpio}{GPIO}{General-Purpose Input/Output}
\newacronym{hal}{HAL}{Hardware Abstraction Layer}
\newacronym{i2c}{I\textsuperscript{2}C}{Inter-Integrated Circuit}
\newacronym{ide}{IDE}{Integrated Development Environment}
\newacronym{ila}{ILA}{Integrated Logic Analyzer}
\newacronym{ism}{ISM}{Industrial, Scientific, and Medical}
\newacronym{isr}{ISR}{Interrupt Service Routine}
\newacronym{lut}{LUT}{Look-Up Table}
\newacronym{mcu}{MCU}{Microcontroller Unit}
\newacronym{ml}{ML}{Machine Learning}
\newacronym{ook}{OOK}{On-Off Keying}
\newacronym{per}{PER}{Packet Error Rate}
\newacronym{rf}{RF}{Radio Frequency}
\newacronym{rssi}{RSSI}{Received Signal Strength Indicator}
\newacronym{rtl}{RTL}{Register-Transfer Level}
\newacronym{snr}{SNR}{Signal-to-Noise Ratio}
\newacronym{spi}{SPI}{Serial Peripheral Interface}
\newacronym{uart}{UART}{Universal Asynchronous Receiver/Transmitter}
\newacronym{vhdl}{VHDL}{VHSIC Hardware Description Language}

% Convenience wrappers so you can type \IDE, \RF, etc.
\providecommand{\ASK}{\gls{ask}}
\providecommand{\BGA}{\gls{bga}}
\providecommand{\BOM}{\gls{bom}}
\providecommand{\BRAM}{\gls{bram}}
\providecommand{\CRC}{\gls{crc}}
\providecommand{\CSS}{\gls{css}}
\providecommand{\DMA}{\gls{dma}}
\providecommand{\DSP}{\gls{dsp}}
\providecommand{\EDA}{\gls{eda}}
\providecommand{\FIFO}{\gls{fifo}}
\providecommand{\FPGA}{\gls{fpga}}
\providecommand{\FSK}{\gls{fsk}}
\providecommand{\FSM}{\gls{fsm}}
\providecommand{\GPIO}{\gls{gpio}}
\providecommand{\HAL}{\gls{hal}}
\providecommand{\IIC}{\gls{i2c}}
\providecommand{\IDE}{\gls{ide}}
\providecommand{\ILA}{\gls{ila}}
\providecommand{\ISM}{\gls{ism}}
\providecommand{\ISR}{\gls{isr}}
\providecommand{\LUT}{\gls{lut}}
\providecommand{\MCU}{\gls{mcu}}
\providecommand{\ML}{\gls{ml}}
\providecommand{\OOK}{\gls{ook}}
\providecommand{\PER}{\gls{per}}
\providecommand{\RF}{\gls{rf}}
\providecommand{\RSSI}{\gls{rssi}}
\providecommand{\RTL}{\gls{rtl}}
\providecommand{\SNR}{\gls{snr}}
\providecommand{\SPI}{\gls{spi}}
\providecommand{\UART}{\gls{uart}}
\providecommand{\VHDL}{\gls{vhdl}}


% Call \printhandoutglossary just before \end{document} in each handout.
\newcommand{\printhandoutglossary}{%
	\clearpage
	\printnoidxglossary[type=\acronymtype,title={Glossary}]%
}

% ── Colors (match Module 1 lesson plan) ──────────────────────
\definecolor{stevens}{RGB}{163,31,52}
\definecolor{darkgray}{RGB}{60,60,60}
\definecolor{lightgray}{RGB}{240,240,240}
\definecolor{accentblue}{RGB}{30,80,140}

\hypersetup{
	colorlinks=true,
	linkcolor=stevens,
	urlcolor=accentblue,
	citecolor=darkgray,
}

% ── Defaults for header variables (handout overrides these) ──
\providecommand{\labid}{?}
\providecommand{\labtitleshort}{Lab Handout}

% ── Running header / footer ──────────────────────────────────
\pagestyle{fancy}
\fancyhf{}
\fancyhead[L]{\small\textcolor{darkgray}{Module 1 --- Lab Handout \labid}}
\fancyhead[R]{\small\textcolor{darkgray}{\labtitleshort}}
\fancyfoot[C]{\small\thepage}
\renewcommand{\headrulewidth}{0.4pt}

% ── Section formatting ───────────────────────────────────────
\titleformat{\section}{\Large\bfseries\color{stevens}}{\thesection}{0.6em}{}
\titleformat{\subsection}{\large\bfseries\color{darkgray}}{\thesubsection}{0.6em}{}

% ── Box styles ───────────────────────────────────────────────
\tcbset{
	infobox/.style={
		enhanced, colback=blue!3, colframe=accentblue, fonttitle=\bfseries,
		title={#1}, coltitle=white, breakable, left=6pt, right=6pt, top=4pt, bottom=4pt
	},
	warning/.style={
		enhanced, colback=red!3, colframe=stevens, fonttitle=\bfseries,
		title={#1}, coltitle=white, breakable, left=6pt, right=6pt, top=4pt, bottom=4pt
	},
	deliverable/.style={
		enhanced, colback=gray!5, colframe=darkgray, fonttitle=\bfseries,
		title={#1}, coltitle=white, breakable, left=6pt, right=6pt, top=4pt, bottom=4pt
	},
}

% ── Figure helper: always draws a placeholder box (images suppressed) ─
% Usage: \labfigure{filename}{width}{caption}{label}
\newcommand{\labfigure}[4]{%
	\begin{figure}[h!]
		\centering
		\fbox{\begin{minipage}[c][2.2in][c]{#2}
			\centering\color{darkgray}\small
				\textbf{[Figure placeholder]}\\[4pt]
				\texttt{\detokenize{#1}}\\[4pt]
				\textit{#3}
		\end{minipage}}%
		\caption{#3}
		\label{#4}
	\end{figure}%
}

% ── Title banner (call at top of document body) ──────────────
\newcommand{\labheader}[2]{%
	\begin{center}
		{\LARGE\bfseries\color{stevens} Lab Handout #1}\\[3pt]
		{\Large\bfseries #2}\\[6pt]
		{\small\color{darkgray} Capstone --- Module 1: \RF{} Hardware Foundations}
	\end{center}
	\vspace{4pt}
	{\color{darkgray}\rule{\textwidth}{0.4pt}}
	\vspace{10pt}
}

% BEFORE \begin{document}. Compiles with a single pdflatex pass (no glossaries tool).

\usepackage[T1]{fontenc}
\usepackage[margin=1in]{geometry}
\usepackage{titlesec}
\usepackage{enumitem}
\usepackage{booktabs}
\usepackage{tabularx}
\usepackage{graphicx}
\graphicspath{{Images/}{../Images/}}
\usepackage{xcolor}
\usepackage{hyperref}
\usepackage{fancyhdr}
\usepackage{amsmath}
\usepackage{amssymb}
\usepackage{tcolorbox}
\usepackage{caption}
\tcbuselibrary{skins,breakable}

% ── Glossary (per-document, built from shared glossary.tex) ──
% Uses \makenoidxglossaries so no external makeglossaries tool is
% needed: sorting and printing is done by LaTeX itself across two
% pdflatex passes (same cadence as the TOC). Only acronyms that
% are actually \gls{}-used in the handout appear in the printed
% glossary, and each \gls{} reference is a live hyperlink that
% jumps to its entry in the Acronyms section at the end.
\usepackage[acronym,nomain,nopostdot]{glossaries}
\makenoidxglossaries
% Shared acronym macro container for EE800/820 lesson plan modules.
% Usage: type \IDE, \FPGA, \UART, etc. First use is expanded; later uses show acronym only.
% Requires in each document preamble:
%   \usepackage[acronym,toc]{glossaries}
%   \makeglossaries
%   % Shared acronym macro container for EE800/820 lesson plan modules.
% Usage: type \IDE, \FPGA, \UART, etc. First use is expanded; later uses show acronym only.
% Requires in each document preamble:
%   \usepackage[acronym,toc]{glossaries}
%   \makeglossaries
%   \input{../glossary}
% Print used acronyms near end of document:
%   \printglossary[type=\acronymtype,title={Acronyms Used}]

% Acronym definitions (only used entries appear in the printed glossary).
\newacronym{ask}{ASK}{Amplitude Shift Keying}
\newacronym{bga}{BGA}{Ball Grid Array}
\newacronym{bom}{BOM}{Bill of Materials}
\newacronym{bram}{BRAM}{Block Random Access Memory}
\newacronym{crc}{CRC}{Cyclic Redundancy Check}
\newacronym{css}{CSS}{Chirp Spread Spectrum}
\newacronym{dma}{DMA}{Direct Memory Access}
\newacronym{dsp}{DSP}{Digital Signal Processing}
\newacronym{eda}{EDA}{Exploratory Data Analysis}
\newacronym{fifo}{FIFO}{First In, First Out}
\newacronym{fpga}{FPGA}{Field-Programmable Gate Array}
\newacronym{fsk}{FSK}{Frequency Shift Keying}
\newacronym{fsm}{FSM}{Finite State Machine}
\newacronym{gpio}{GPIO}{General-Purpose Input/Output}
\newacronym{hal}{HAL}{Hardware Abstraction Layer}
\newacronym{i2c}{I\textsuperscript{2}C}{Inter-Integrated Circuit}
\newacronym{ide}{IDE}{Integrated Development Environment}
\newacronym{ila}{ILA}{Integrated Logic Analyzer}
\newacronym{ism}{ISM}{Industrial, Scientific, and Medical}
\newacronym{isr}{ISR}{Interrupt Service Routine}
\newacronym{lut}{LUT}{Look-Up Table}
\newacronym{mcu}{MCU}{Microcontroller Unit}
\newacronym{ml}{ML}{Machine Learning}
\newacronym{ook}{OOK}{On-Off Keying}
\newacronym{per}{PER}{Packet Error Rate}
\newacronym{rf}{RF}{Radio Frequency}
\newacronym{rssi}{RSSI}{Received Signal Strength Indicator}
\newacronym{rtl}{RTL}{Register-Transfer Level}
\newacronym{snr}{SNR}{Signal-to-Noise Ratio}
\newacronym{spi}{SPI}{Serial Peripheral Interface}
\newacronym{uart}{UART}{Universal Asynchronous Receiver/Transmitter}
\newacronym{vhdl}{VHDL}{VHSIC Hardware Description Language}

% Convenience wrappers so you can type \IDE, \RF, etc.
\providecommand{\ASK}{\gls{ask}}
\providecommand{\BGA}{\gls{bga}}
\providecommand{\BOM}{\gls{bom}}
\providecommand{\BRAM}{\gls{bram}}
\providecommand{\CRC}{\gls{crc}}
\providecommand{\CSS}{\gls{css}}
\providecommand{\DMA}{\gls{dma}}
\providecommand{\DSP}{\gls{dsp}}
\providecommand{\EDA}{\gls{eda}}
\providecommand{\FIFO}{\gls{fifo}}
\providecommand{\FPGA}{\gls{fpga}}
\providecommand{\FSK}{\gls{fsk}}
\providecommand{\FSM}{\gls{fsm}}
\providecommand{\GPIO}{\gls{gpio}}
\providecommand{\HAL}{\gls{hal}}
\providecommand{\IIC}{\gls{i2c}}
\providecommand{\IDE}{\gls{ide}}
\providecommand{\ILA}{\gls{ila}}
\providecommand{\ISM}{\gls{ism}}
\providecommand{\ISR}{\gls{isr}}
\providecommand{\LUT}{\gls{lut}}
\providecommand{\MCU}{\gls{mcu}}
\providecommand{\ML}{\gls{ml}}
\providecommand{\OOK}{\gls{ook}}
\providecommand{\PER}{\gls{per}}
\providecommand{\RF}{\gls{rf}}
\providecommand{\RSSI}{\gls{rssi}}
\providecommand{\RTL}{\gls{rtl}}
\providecommand{\SNR}{\gls{snr}}
\providecommand{\SPI}{\gls{spi}}
\providecommand{\UART}{\gls{uart}}
\providecommand{\VHDL}{\gls{vhdl}}


% Print used acronyms near end of document:
%   \printglossary[type=\acronymtype,title={Acronyms Used}]

% Acronym definitions (only used entries appear in the printed glossary).
\newacronym{ask}{ASK}{Amplitude Shift Keying}
\newacronym{bga}{BGA}{Ball Grid Array}
\newacronym{bom}{BOM}{Bill of Materials}
\newacronym{bram}{BRAM}{Block Random Access Memory}
\newacronym{crc}{CRC}{Cyclic Redundancy Check}
\newacronym{css}{CSS}{Chirp Spread Spectrum}
\newacronym{dma}{DMA}{Direct Memory Access}
\newacronym{dsp}{DSP}{Digital Signal Processing}
\newacronym{eda}{EDA}{Exploratory Data Analysis}
\newacronym{fifo}{FIFO}{First In, First Out}
\newacronym{fpga}{FPGA}{Field-Programmable Gate Array}
\newacronym{fsk}{FSK}{Frequency Shift Keying}
\newacronym{fsm}{FSM}{Finite State Machine}
\newacronym{gpio}{GPIO}{General-Purpose Input/Output}
\newacronym{hal}{HAL}{Hardware Abstraction Layer}
\newacronym{i2c}{I\textsuperscript{2}C}{Inter-Integrated Circuit}
\newacronym{ide}{IDE}{Integrated Development Environment}
\newacronym{ila}{ILA}{Integrated Logic Analyzer}
\newacronym{ism}{ISM}{Industrial, Scientific, and Medical}
\newacronym{isr}{ISR}{Interrupt Service Routine}
\newacronym{lut}{LUT}{Look-Up Table}
\newacronym{mcu}{MCU}{Microcontroller Unit}
\newacronym{ml}{ML}{Machine Learning}
\newacronym{ook}{OOK}{On-Off Keying}
\newacronym{per}{PER}{Packet Error Rate}
\newacronym{rf}{RF}{Radio Frequency}
\newacronym{rssi}{RSSI}{Received Signal Strength Indicator}
\newacronym{rtl}{RTL}{Register-Transfer Level}
\newacronym{snr}{SNR}{Signal-to-Noise Ratio}
\newacronym{spi}{SPI}{Serial Peripheral Interface}
\newacronym{uart}{UART}{Universal Asynchronous Receiver/Transmitter}
\newacronym{vhdl}{VHDL}{VHSIC Hardware Description Language}

% Convenience wrappers so you can type \IDE, \RF, etc.
\providecommand{\ASK}{\gls{ask}}
\providecommand{\BGA}{\gls{bga}}
\providecommand{\BOM}{\gls{bom}}
\providecommand{\BRAM}{\gls{bram}}
\providecommand{\CRC}{\gls{crc}}
\providecommand{\CSS}{\gls{css}}
\providecommand{\DMA}{\gls{dma}}
\providecommand{\DSP}{\gls{dsp}}
\providecommand{\EDA}{\gls{eda}}
\providecommand{\FIFO}{\gls{fifo}}
\providecommand{\FPGA}{\gls{fpga}}
\providecommand{\FSK}{\gls{fsk}}
\providecommand{\FSM}{\gls{fsm}}
\providecommand{\GPIO}{\gls{gpio}}
\providecommand{\HAL}{\gls{hal}}
\providecommand{\IIC}{\gls{i2c}}
\providecommand{\IDE}{\gls{ide}}
\providecommand{\ILA}{\gls{ila}}
\providecommand{\ISM}{\gls{ism}}
\providecommand{\ISR}{\gls{isr}}
\providecommand{\LUT}{\gls{lut}}
\providecommand{\MCU}{\gls{mcu}}
\providecommand{\ML}{\gls{ml}}
\providecommand{\OOK}{\gls{ook}}
\providecommand{\PER}{\gls{per}}
\providecommand{\RF}{\gls{rf}}
\providecommand{\RSSI}{\gls{rssi}}
\providecommand{\RTL}{\gls{rtl}}
\providecommand{\SNR}{\gls{snr}}
\providecommand{\SPI}{\gls{spi}}
\providecommand{\UART}{\gls{uart}}
\providecommand{\VHDL}{\gls{vhdl}}


% Call \printhandoutglossary just before \end{document} in each handout.
\newcommand{\printhandoutglossary}{%
	\clearpage
	\printnoidxglossary[type=\acronymtype,title={Glossary}]%
}

% ── Colors (match Module 1 lesson plan) ──────────────────────
\definecolor{stevens}{RGB}{163,31,52}
\definecolor{darkgray}{RGB}{60,60,60}
\definecolor{lightgray}{RGB}{240,240,240}
\definecolor{accentblue}{RGB}{30,80,140}

\hypersetup{
	colorlinks=true,
	linkcolor=stevens,
	urlcolor=accentblue,
	citecolor=darkgray,
}

% ── Defaults for header variables (handout overrides these) ──
\providecommand{\labid}{?}
\providecommand{\labtitleshort}{Lab Handout}

% ── Running header / footer ──────────────────────────────────
\pagestyle{fancy}
\fancyhf{}
\fancyhead[L]{\small\textcolor{darkgray}{Module 1 --- Lab Handout \labid}}
\fancyhead[R]{\small\textcolor{darkgray}{\labtitleshort}}
\fancyfoot[C]{\small\thepage}
\renewcommand{\headrulewidth}{0.4pt}

% ── Section formatting ───────────────────────────────────────
\titleformat{\section}{\Large\bfseries\color{stevens}}{\thesection}{0.6em}{}
\titleformat{\subsection}{\large\bfseries\color{darkgray}}{\thesubsection}{0.6em}{}

% ── Box styles ───────────────────────────────────────────────
\tcbset{
	infobox/.style={
		enhanced, colback=blue!3, colframe=accentblue, fonttitle=\bfseries,
		title={#1}, coltitle=white, breakable, left=6pt, right=6pt, top=4pt, bottom=4pt
	},
	warning/.style={
		enhanced, colback=red!3, colframe=stevens, fonttitle=\bfseries,
		title={#1}, coltitle=white, breakable, left=6pt, right=6pt, top=4pt, bottom=4pt
	},
	deliverable/.style={
		enhanced, colback=gray!5, colframe=darkgray, fonttitle=\bfseries,
		title={#1}, coltitle=white, breakable, left=6pt, right=6pt, top=4pt, bottom=4pt
	},
}

% ── Figure helper: always draws a placeholder box (images suppressed) ─
% Usage: \labfigure{filename}{width}{caption}{label}
\newcommand{\labfigure}[4]{%
	\begin{figure}[h!]
		\centering
		\fbox{\begin{minipage}[c][2.2in][c]{#2}
			\centering\color{darkgray}\small
				\textbf{[Figure placeholder]}\\[4pt]
				\texttt{\detokenize{#1}}\\[4pt]
				\textit{#3}
		\end{minipage}}%
		\caption{#3}
		\label{#4}
	\end{figure}%
}

% ── Title banner (call at top of document body) ──────────────
\newcommand{\labheader}[2]{%
	\begin{center}
		{\LARGE\bfseries\color{stevens} Lab Handout #1}\\[3pt]
		{\Large\bfseries #2}\\[6pt]
		{\small\color{darkgray} Capstone --- Module 1: \RF{} Hardware Foundations}
	\end{center}
	\vspace{4pt}
	{\color{darkgray}\rule{\textwidth}{0.4pt}}
	\vspace{10pt}
}


\begin{document}
\labheader{1-G}{Receiver Firmware: ISR-Driven Capture and \UART{} Forwarding}

\section{Objective}
Implement receiver firmware that captures beacon packets via an \texttt{RxDone} interrupt, queues them in a ring buffer, and forwards them from the main loop to the \FPGA{} over a 43-byte framed \UART{} protocol. This is the software side of the Week~4 deliverable.

\section{Prerequisites}
\begin{itemize}[leftmargin=2em]
	\item Lab Handouts~1-A, 1-B, 1-C, 1-D, 1-F complete.
	\item At least one working beacon from LH1-F as a traffic source.
\end{itemize}

\section{Forwarding Frame Format}
The receiver wraps each received beacon in a framed \UART{} message for the \FPGA{}. Total 43 bytes:
\begin{table}[h!]
	\centering
	\begin{tabular}{c l l}
		\toprule
		\textbf{Offset} & \textbf{Field} & \textbf{Length} \\
		\midrule
		0 & Start delimiter (\texttt{0x7E}) & 1 \\
		1 & Frame length (always 41) & 1 \\
		2--33 & Original 32-byte beacon packet & 32 \\
		34--35 & \RSSI{} (int16, \mbox{LE}, dBm $\times$1) & 2 \\
		36 & \SNR{} (int8, dB $\times$4 per SX1276 scaling) & 1 \\
		37--40 & Receiver timestamp (uint32, \mbox{LE}, ms) & 4 \\
		41--42 & \CRC{}-16/CCITT-FALSE over bytes 0--40 & 2 \\
		\bottomrule
	\end{tabular}
\end{table}

\section{Architecture}

\subsection{ISR path}
\begin{enumerate}[leftmargin=2em]
	\item DIO0 rising edge triggers an \texttt{EXTI} interrupt.
	\item \ISR{} reads \texttt{RegIrqFlags}. If \texttt{RxDone} is set and \texttt{PayloadCrcError} is clear:
	\begin{enumerate}[leftmargin=1.2em, label=(\alph*)]
		\item Read \texttt{RegFifoRxCurrentAddr}, set \texttt{RegFifoAddrPtr}, read \texttt{RegRxNbBytes} bytes from \texttt{RegFifo}.
		\item Read \texttt{RegPktRssiValue} and \texttt{RegPktSnrValue}.
		\item Stamp the current \texttt{ms\_count}.
		\item Push a struct (\texttt{beacon\_pkt\_t}) containing these fields into a \textit{lock-free single-producer/single-consumer ring buffer} (producer = \ISR{}, consumer = main loop).
		\item Clear \texttt{RegIrqFlags} by writing \texttt{0xFF}.
	\end{enumerate}
\end{enumerate}

\subsection{Main-loop path}
\begin{enumerate}[leftmargin=2em]
	\item Check ring buffer head vs.\ tail. If not empty, dequeue one packet.
	\item Build the 43-byte forwarding frame.
	\item Compute \CRC{}-16/CCITT-FALSE over bytes 0--40 and place in bytes 41--42.
	\item Transmit the entire 43-byte frame over \texttt{USART3} (or whichever \UART{} is wired to the Nexys A7) using \texttt{HAL\_UART\_Transmit\_DMA}.
	\item Increment a \texttt{frames\_forwarded} counter.
\end{enumerate}

\subsection{Ring buffer sizing}
A 16-entry ring ($\approx$640 bytes of beacon structs) accommodates roughly 16 back-to-back receptions without overflow. With two beacons at 2 and 3 second intervals, this is $\sim$40\,s of margin against a stalled main loop --- comfortable.

\section{Procedure}

\subsection{Step 1 --- Peripheral configuration}
\begin{enumerate}[leftmargin=2em]
	\item Starting from your LH1-F project or a new one, configure:
	\begin{itemize}[leftmargin=1.2em]
		\item \texttt{SPI1} (unchanged from LH1-C).
		\item \texttt{EXTI} on \texttt{PA8} (DIO0), rising-edge trigger.
		\item \texttt{USART3} at 115200\,8N1 on \texttt{PC4}/\texttt{PC5} (TX/RX), with TX \mbox{DMA} enabled.
		\item \texttt{USART2} still enabled for \texttt{printf} debug over the ST-LINK virtual COM port.
		\item \texttt{TIM2} 1\,ms tick (unchanged).
	\end{itemize}
\end{enumerate}

\subsection{Step 2 --- Ring buffer implementation}
\begin{enumerate}[leftmargin=2em]
	\item Define the packet struct:
\begin{verbatim}
typedef struct {
    uint8_t  raw[32];
    int16_t  rssi_dbm;
    int8_t   snr_q4;
    uint32_t rx_timestamp_ms;
} beacon_pkt_t;
\end{verbatim}
	\item Static ring with power-of-two size: \texttt{static beacon\_pkt\_t ring[16];} and \texttt{volatile uint32\_t head, tail;}.
	\item \texttt{push}: producer (\ISR{}). Fails silently on overflow and increments a \texttt{dropped} counter.
	\item \texttt{pop}: consumer (main loop). Returns bool indicating whether a packet was dequeued.
	\item Use a single \texttt{\_\_DMB()} memory barrier between buffer write and head increment to ensure ordering on Cortex-M4.
\end{enumerate}

\subsection{Step 3 --- ISR}
\begin{enumerate}[leftmargin=2em]
	\item Implement \texttt{HAL\_GPIO\_EXTI\_Callback(uint16\_t GPIO\_Pin)}. Check that the pin matches DIO0.
	\item Keep the \ISR{} short: FIFO read, \RSSI{}/\SNR{} read, timestamp, \texttt{push} to ring, clear IRQ flags. No \texttt{printf} inside the \ISR{}.
\end{enumerate}

\subsection{Step 4 --- Frame builder and forwarder}
\begin{enumerate}[leftmargin=2em]
	\item In the main loop, call \texttt{pop}. If a packet is present, construct the 43-byte frame as described above.
	\item Reuse the \CRC{}-16 helper from LH1-F.
	\item Call \texttt{HAL\_UART\_Transmit\_DMA(\&huart3, frame, 43)}. Track the \mbox{DMA} completion flag; do not overwrite the frame buffer until transmission is complete.
	\item Also \texttt{printf} a human-readable summary over USART2 for debug: \texttt{"RX beacon=\%02X seq=\%lu rssi=\%d snr=\%d"}.
\end{enumerate}

\subsection{Step 5 --- Validation without the \FPGA{}}
\begin{enumerate}[leftmargin=2em]
	\item Before wiring to the Nexys A7, loopback-test the forwarding \UART{} by connecting \texttt{USART3 TX} to \texttt{USART3 RX} (or to a USB-\UART{} adapter) and capturing frames on a host PC.
	\item Use a Python script (\texttt{pyserial}) to parse incoming 43-byte frames, verify the start delimiter, validate the \CRC{}-16, and print the extracted beacon \mbox{ID}, sequence number, \RSSI{}, and \SNR{}.
	\item Run with Beacon~A (and optionally Beacon~B) from LH1-F broadcasting.
\end{enumerate}

\section{Verification Checklist}
\begin{itemize}[leftmargin=2em]
	\item[$\square$] \texttt{RxDone} interrupt fires on every received packet (\ISR{} entry count matches expected receptions).
	\item[$\square$] Ring buffer overflows only if you deliberately stall the main loop.
	\item[$\square$] 43-byte frames appear on \texttt{USART3 TX} for every received beacon.
	\item[$\square$] Host Python parser validates start delimiter and \CRC{}-16 on 100\% of frames.
	\item[$\square$] Dual-beacon run: parser sees interleaved \mbox{ID}\,=\,\texttt{0x01} and \mbox{ID}\,=\,\texttt{0x02} frames with independently contiguous sequence numbers.
	\item[$\square$] Dropped-packet counter remains zero during a 5-minute run.
\end{itemize}

\begin{tcolorbox}[warning={Concurrency Hazards}]
	\begin{itemize}[leftmargin=1.2em]
		\item Access \texttt{head}/\texttt{tail} as \texttt{volatile}. Never read the same index twice in the consumer without rechecking.
		\item Do not call \HAL{} \texttt{printf} from inside the \ISR{} --- it blocks on \mbox{UART} transmit and will drop \RF{} packets.
		\item Do not overwrite the \mbox{DMA} source buffer while a prior transmission is still in flight. Use a simple busy flag or a second buffer.
	\end{itemize}
\end{tcolorbox}

\begin{tcolorbox}[deliverable={Lab Handout 1-G Deliverable}]
	\begin{enumerate}[leftmargin=1.2em]
		\item Full receiver firmware source, including ring buffer and \ISR{}.
		\item Python parser script (one file) that validates frames from \texttt{USART3}.
		\item 5-minute dual-beacon log showing zero dropped packets and zero \CRC{} failures in the parser.
		\item One-paragraph rationale for the ring buffer size you chose, including a worst-case arrival-rate analysis.
	\end{enumerate}
\end{tcolorbox}

\end{document}
